\section{Systems and Experimental Design}
\label{sec:systems}
Table~\ref{tab:systems} defines the five arms. For a given LLM, generation and selection use the same configured model. Across LLMs, the prompt text is shared within each task, and selection uses the same displayed candidate order for each item. The implemented order is a seeded shuffle of the stored lists; it is recorded with each response. These controls align inputs but do not make architectures, training or inference equivalent.

\begin{table}[t]
\centering\small
\begin{tabular}{@{}lll@{}}
\toprule
System & Candidates & Output \\
\midrule
DictaBERT ranker & Yes & Top candidate \\
Qwen generation & No & Expansion text \\
Qwen selection & Yes & Candidate letter \\
Gemini generation & No & Expansion text \\
Gemini selection & Yes & Candidate letter \\
\bottomrule
\end{tabular}
\caption{All arms receive a sentence and marked target. Only DictaBERT receives task-specific supervised training. Saved development runs cover these configurations, with partial Gemini execution.}
\label{tab:systems}
\end{table}

\subsection{Supervised Hebrew candidate ranker}
DictaBERT is a pretrained Hebrew encoder \citep{shmidman2023dictabert}. Our implementation is a \emph{cross-encoder}: the marked sentence and one candidate are encoded jointly, allowing attention between their tokens. A dropout layer and a linear head turn the pooled representation into a scalar score. The reconstructed Colab model uses the initial classification token (CLS). Candidates are scored independently, and prediction chooses the largest score within the item's inventory.

For each training occurrence, the reference candidate is positive and the other supplied candidates are negative. The implemented objective is binary cross-entropy on the candidate-pair logits:
\begin{equation}
\begin{split}
\mathcal{L}=-\frac{1}{M}\sum_{j=1}^{M}\big[&z_j\log\sigma(u_j)\,+\\
&(1-z_j)\log(1-\sigma(u_j))\big],
\end{split}
\end{equation}
where $M$ is the number of pairs, $u_j$ is a score, $\sigma$ is the sigmoid function and $z_j\in\{0,1\}$ is the label. Thus occurrences with larger inventories contribute more pairs. This is binary pair training, not a multiclass softmax over each inventory. AdamW updates the encoder and scoring head; no parameter-efficient adapter is used in this implementation. Full fine-tuning is a change from the head-or-adapter options in the submitted proposal.

The inspected notebook starts from \texttt{dicta-il/dictabert} and specifies CLS pooling, dropout 0.1, learning rate $2\times10^{-5}$, pair batch size 16, one epoch and seed 42. These are inspected notebook settings, not a recovered execution log. It sets the seed \emph{after} initializing the model and added embeddings, so this seed does not control their initialization. It retains a checkpoint on strict improvement in development pair loss and saves the complete \texttt{state\_dict} in fp16. Pair loss differs from occurrence-level selection accuracy. The separate package training route sets the seed before initialization and requires a manifest; it is not the evidenced source of these weights.

The pair layout is a marked context followed by the candidate, separated by BERT special tokens. The nominal length limit is 256. Context is cropped from the outside while retaining the target, markers and candidate. The Colab routine can exceed that limit if the required tokens alone do not fit; the package predictor instead rejects that case. Saved checks show identical encoding tensors for all 62 dev items and matching forward logits on three validation items. Candidate length can affect the context budget, whereas the LLM receives the full sentence.

Shaked confirmed that the supplied weights came from Ben's current Colab run. The saved reconstruction verifies their SHA-256 and a strict load of all 201 tensors, including the encoder, two added marker embeddings and scoring head, with no missing or unexpected keys and no random-head fallback. The saved tensors are fp16; development inference used fp32 on CPU. The record also identifies the local tokenizer and configuration. This supports reconstruction and prediction, not complete reproduction of training. The exact original input snapshot, resolved base-model/tokenizer revision, library versions, hardware, selected epoch and loss remain undocumented. The filename suggests seed 43 while the notebook specifies 42; neither establishes the executed seed.

\subsection{Prompted language models}
The generation prompt, written in Hebrew, asks for the expansion of the marked acronym in a few words, without explanation. The selection prompt supplies lettered candidates and requests exactly one letter. Neither contains demonstrations or identifies the correct answer. Full prompt templates belong to the technical companion; their information content and response requirements are the experimental distinction described here.

The saved Qwen runs use \texttt{qwen2.5:7b}, from the Qwen2.5 family \citep{qwen2025technical}, through Ollama. Runtime metadata reports Q4\_K\_M quantization and the same server model digest across all 124 generation and selection responses. No additional decoding options were supplied, so the run is not described as deterministic. Gemini requests use \texttt{gemini-3.8-flash} with thinking level \texttt{low}; successful responses report the same version string. A provider version string is weaker identity evidence than a content digest, and low thinking does not mean disabled thinking. Omitted service defaults remain unresolved.

Requests use a 120-second timeout. The basic request path has no automatic retries; the Gemini continuation used at most three attempts per item for transient errors, with backoff and saved progress. Completed responses were retained across continuation. Artifacts preserve prompts, candidate orders, raw responses, completion status, model identity and failures. Gemini collection remained partial after HTTP 429 and 503 errors. Service availability is reported separately from answer quality, without attributing the failures to an unverified quota mechanism.
